\chapter{Spis skrótów i symboli}

\begin{itemize}
    \item ATM – Automatyczna Transkrypcja Muzyczna (\english{Automatic Music Transcription})
    \item GTT – Transkrypcja Gitarowa na Tabulaturę (\english{Guitar Tablature Transcription})
    \item MIR – Komputerowe Pozyskiwanie Informacji Muzycznej (\english{Music Information Retrieval})
    \item CPS – Cyfrowe Przetwarzanie Sygnałów (\english{Digital Signal Processing – DSP})
    \item UM – Uczenie Maszynowe (\english{Machine Learning – ML})
    \item SI – Sztuczna Inteligencja (\english{Artificial Intelligence – AI})
    \item DFT – Dyskretna Transformata Fouriera (\english{Discrete Fourier Transform})
    \item STFT – Krótkoczasowa Transformata Fouriera (\english{Short-Time Fourier Transform})
    \item CQT – Transformata \english{Constant-Q} (\english{Constant-Q Transform})
    \item MPE – Estymacja Wielokrotnej Wysokości Dźwięku (\english{Multiple Pitch Estimation})
    \item NMF – Nieujemna Faktoryzacja Macierzy (\english{Non-negative Matrix Factorization})
    \item HMM – Ukryte Modele Markowa (\english{Hidden Markov Models})
    \item PLCA – Probabilistyczna Analiza Składowych Ukrytych (\english{Probabilistic Latent Component Analysis})
    \item DL – Głębokie Uczenie (\english{Deep Learning})
    \item CNN – Splotowe Sieci Neuronowe (\english{Convolutional Neural Networks})
    \item RNN – Rekurencyjne Sieci Neuronowe (\english{Recurrent Neural Networks})
    \item LSTM – Długa Pamięć Krótkoterminowa (\english{Long Short-Term Memory})
    \item GRU – Bramkowana Jednostka Rekurencyjna (\english{Gated Recurrent Unit})
    \item CRNN – Konwolucyjno-Rekurencyjne Sieci Neuronowe (\english{Convolutional Recurrent Neural Networks})
    \item $F_0$ – Częstotliwość podstawowa
    \item $X_k$ – k-ty współczynnik dyskretnej transformaty Fouriera
    \item $x_n$ – n-ta próbka sygnału w dziedzinie czasu
    \item $N$ – długość ciągu próbek sygnału (w kontekście DFT)
    \item $i$ – jednostka urojona
\end{itemize}